\section{Mapeamento de canais}
\label{sec:mapeamento}

Este capítulo lista os canais que o \projeto{} pode utilizar para alcançar seus clientes,
organizados nas três categorias solicitadas: comunicação, distribuição/entrega e
relacionamento. O mapeamento parte do Modelo de Negócio (Atividade~III), da Proposta de Valor
e Segmentação (Atividade~IV) e das entrevistas com tutores realizadas nas etapas anteriores.

Como o \projeto{} opera uma plataforma com múltiplos lados, os canais precisam alcançar quatro
públicos: o \textbf{tutor de baixa renda} (classes C, D e E), o \textbf{tutor voluntário} (que
acumula créditos com serviços comunitários), o \textbf{prestador veterinário parceiro} e as
\textbf{ONGs e protetores}. Os canais marcados com \prio{} foram selecionados como
prioritários no Capítulo~\ref{sec:estrategia}.

% ---------------------------------------------------------------------
\subsection{Canais de comunicação}
\label{sec:canais-comunicacao}

Definem como o cliente fica sabendo da startup.

\begin{longtable}{L{3.7cm} L{7.0cm} L{3.2cm}}
\toprule
\cabfundo \cab{Canal} & \cab{Descrição} & \cab{Público alcançado} \\
\midrule
\endfirsthead
\toprule
\cabfundo \cab{Canal} & \cab{Descrição} & \cab{Público alcançado} \\
\midrule
\endhead
\bottomrule
\endfoot
\prio{} \canal{Boca a boca comunitário e ONGs} & Divulgação por ONGs, protetores independentes, grupos de WhatsApp de bairro e lideranças locais, apoiada em confiança prévia. & Tutores, voluntários, ONGs \\
\canal{Programa de indicação} & Tutor convida outro tutor e ambos recebem bônus de créditos ao concluir o cadastro do pet. & Tutores, voluntários \\
\prio{} \canal{Pontos de atendimento parceiros} & QR code e material impresso em clínicas, pet shops e agropecuárias, captando o tutor no momento da necessidade. & Tutores de baixa renda \\
\prio{} \canal{Eventos presenciais} & Campanhas de vacinação, feiras de adoção e ações do poder público, com cadastro assistido do pet. & Tutores, ONGs, poder público \\
\canal{Redes sociais orgânicas} & Conteúdo educativo de prevenção (vacinas, vermifugação, sinais de alerta) em Instagram, TikTok e Facebook. & Tutores, voluntários \\
\canal{Parcerias com o poder público} & Divulgação com centros de controle de zoonoses, UBS Animal e secretarias de saúde. & Tutores de baixa renda \\
\canal{Faculdades de Medicina Veterinária} & Ações de extensão, ligas acadêmicas e atendimentos supervisionados que divulgam a plataforma. & Prestadores, voluntários \\
\canal{Mídia local e rádios comunitárias} & Inserções e entrevistas em veículos de bairro e rádios comunitárias. & Tutores de baixa renda \\
\canal{Influenciadores pet locais} & Microinfluenciadores e criadores de conteúdo pet da região, por permuta. & Tutores, voluntários \\
\canal{Mídia paga digital} & Anúncios geolocalizados (Meta Ads e Google Ads) para instalação do aplicativo. & Tutores, voluntários \\
\end{longtable}

% ---------------------------------------------------------------------
\subsection{Canais de distribuição/entrega}
\label{sec:canais-distribuicao}

Definem como o cliente recebe a solução, tanto o aplicativo quanto o atendimento veterinário
que ele viabiliza.

\begin{longtable}{L{3.7cm} L{7.0cm} L{3.2cm}}
\toprule
\cabfundo \cab{Canal} & \cab{Descrição} & \cab{Público alcançado} \\
\midrule
\endfirsthead
\toprule
\cabfundo \cab{Canal} & \cab{Descrição} & \cab{Público alcançado} \\
\midrule
\endhead
\bottomrule
\endfoot
\prio{} \canal{Google Play Store} & Distribuição do aplicativo Android; a primeira fase é 100\% mobile, coerente com o acesso do público-alvo. & Todos os segmentos \\
\prio{} \canal{Rede de clínicas e veterinários parceiros} & Entrega presencial do atendimento (consulta, vacina, exame) com preço social e resgate de créditos, agendado pelo aplicativo. & Tutores, prestadores \\
\prio{} \canal{Cadastro e instalação assistidos} & Voluntário ou equipe instala o app e cadastra o pet junto ao tutor em eventos e pontos parceiros. & Tutores com menor familiaridade digital \\
\canal{Mutirões de vacinação e castração} & Entrega do cuidado preventivo em campanhas coletivas com parceiros públicos e ONGs. & Tutores, ONGs \\
\canal{Prontuário digital via QR Code} & Acesso do profissional ao histórico do pet durante o atendimento, sem cadastro adicional. & Tutores, prestadores \\
\canal{Gateway de pagamento (terceiros)} & Adquirente que viabiliza o \emph{check-out} híbrido (créditos + dinheiro) e o parcelamento. & Tutores, prestadores \\
\canal{App Store (iOS)} & Publicação para iPhone em fase posterior, aproveitando a base de código multiplataforma. & Tutores, voluntários \\
\canal{Aplicação web (PWA)} & Versão acessível pelo navegador; \emph{descartada na primeira fase}, conforme a Proposta de Solução. & Prestadores (futuro) \\
\end{longtable}

% ---------------------------------------------------------------------
\subsection{Canais de relacionamento}
\label{sec:canais-relacionamento}

Definem como a startup mantém contato com o cliente após a aquisição.

\begin{longtable}{L{3.7cm} L{7.0cm} L{3.2cm}}
\toprule
\cabfundo \cab{Canal} & \cab{Descrição} & \cab{Público alcançado} \\
\midrule
\endfirsthead
\toprule
\cabfundo \cab{Canal} & \cab{Descrição} & \cab{Público alcançado} \\
\midrule
\endhead
\bottomrule
\endfoot
\prio{} \canal{Notificações push automatizadas} & Lembretes de vacinação, vermifugação e retorno gerados a partir do calendário do pet; principal mecanismo de retenção. & Tutores \\
\prio{} \canal{Comunidade de tutores e banco de horas} & Troca de serviços entre vizinhos com reputação e avaliação mútuas, gerando pertencimento. & Tutores, voluntários \\
\prio{} \canal{Suporte humano assistido} & Atendimento por mensagens (WhatsApp Business e chat no app) para casos sensíveis, disputas de crédito e orientação. & Tutores, voluntários \\
\canal{Autosserviço guiado} & \emph{Onboarding} do pet em poucos passos, que já gera o calendário preventivo. & Tutores \\
\canal{Triagem orientativa} & Chat estruturado e não diagnóstico que indica o nível de urgência de um sintoma. & Tutores \\
\canal{Gestor de relacionamento com prestadores} & Acompanhamento de volume, repasses e satisfação das clínicas parceiras. & Prestadores \\
\canal{Campanhas coletivas} & Mutirões e ações sazonais como momento de reengajamento da base. & Tutores, ONGs \\
\canal{Pesquisas de satisfação no app} & NPS e CSAT após atendimentos e trocas de créditos. & Todos os segmentos \\
\end{longtable}
